%%=========================================================================
%%  Section II: System Model
%%  WCSP Conference Paper — IEEE Two-Column Format
%%=========================================================================

\section{System Model}

We consider a LEO satellite non-terrestrial network (NTN) downlink operating in the S-band. The satellite transmits an OTFS-modulated signal carrying both communication data and an embedded delay-Doppler (DD) domain pilot for integrated sensing and communication (ISAC). Fig.~\ref{fig:system_diagram} illustrates the overall system architecture.

%% -----------------------------------------------------------------------
\subsection{OTFS Modulation and DD-Domain Signal Representation}

The OTFS system occupies a bandwidth of $B = N \Delta f$ using $N$ subcarriers with subcarrier spacing $\Delta f$, and spans $M$ OFDM symbols per frame. The DD grid has $N \times M$ resource elements with delay resolution $\Delta\tau = 1/(N\Delta f)$ and Doppler resolution $\Delta\nu = 1/(M T_{\mathrm{s}})$, where $T_{\mathrm{s}} = (1+\beta)/\Delta f$ is the OFDM symbol duration including a cyclic prefix (CP) ratio $\beta$.

At the transmitter, $N_{\mathrm{d}}$ QAM data symbols and one pilot symbol are mapped onto the DD grid $\mathbf{X}^{\mathrm{DD}} \in \mathbb{C}^{N \times M}$. The DD grid is transformed to the time-frequency (TF) domain via the Inverse Symplectic Finite Fourier Transform (ISFFT):
\begin{equation}\label{eq:isfft}
  X^{\mathrm{TF}}[n,m] = \frac{1}{\sqrt{NM}} \sum_{l=0}^{N-1} \sum_{k=0}^{M-1} X^{\mathrm{DD}}[l,k]\, e^{j2\pi\left(\frac{nl}{N} - \frac{mk}{M}\right)},
\end{equation}
where $l$ and $k$ denote the delay and Doppler bin indices, and $n$ and $m$ denote the subcarrier and OFDM symbol indices, respectively. The TF-domain signal is then transmitted using standard OFDM processing with CP insertion.

At the receiver, after CP removal and FFT demodulation, the received TF grid is converted back to the DD domain via the Symplectic Finite Fourier Transform (SFFT):
\begin{equation}\label{eq:sfft}
  Y^{\mathrm{DD}}[l,k] = \frac{1}{\sqrt{NM}} \sum_{n=0}^{N-1} \sum_{m=0}^{M-1} Y^{\mathrm{TF}}[n,m]\, e^{-j2\pi\left(\frac{nl}{N} - \frac{mk}{M}\right)}.
\end{equation}

%% -----------------------------------------------------------------------
\subsection{LEO NTN Channel Model}

The LEO satellite-to-ground channel follows the 3GPP TR~38.811 NTN-TDL model~\cite{3gpp_38811}. The channel consists of $P$ propagation paths, each characterized by a delay $\tau_i$, Doppler shift $\nu_i$, and complex gain $h_i$, for $i = 1, \ldots, P$.

\subsubsection{Multipath Structure}
We adopt the NTN-TDL-C profile, which is derived from the NTN-TDL-A NLOS model via Rician K-factor scaling~\cite{3gpp_38901}. The base tap delays are $\tilde{\tau}_i = d_i \cdot \tau_{\mathrm{ref}}$, where $d_i \in \{0,\; 1.08,\; 2.84\}$ are the normalized delays and $\tau_{\mathrm{ref}} = 30$~ns is the reference delay spread. For the LOS profile, the path powers are scaled as:
\begin{equation}\label{eq:kfactor_scaling}
  P_i^{\mathrm{scatter}} = \frac{P_i^{\mathrm{base}}}{K + 1}, \quad
  P^{\mathrm{LoS}} = \frac{K}{K + 1},
\end{equation}
where $K$ is the Rician K-factor in linear scale, determined by the satellite elevation angle $\theta_{\mathrm{el}}$ and the propagation scenario (e.g., Suburban, Urban). For example, at $\theta_{\mathrm{el}} = 50^\circ$ in a Suburban environment, $K \approx 12$~dB~\cite{3gpp_38811}.

\subsubsection{Shadowed-Rician LoS Fading}
The LoS component undergoes shadowed-Rician (Loo model) fading to capture ionospheric scintillation and tropospheric shadowing~\cite{loo_model}. The LoS amplitude is modulated by a lognormal shadow factor:
\begin{equation}\label{eq:los_gain}
  h_1 = \sqrt{\frac{K}{K+1}} \cdot 10^{\xi/20} \cdot e^{j\phi_0} + h_1^{\mathrm{scatter}},
\end{equation}
where $\xi \sim \mathcal{N}(0, \sigma_{\mathrm{s}}^2)$ is the shadow fading in dB, $\phi_0 \sim \mathcal{U}[0, 2\pi)$ is a random phase, and $\sigma_{\mathrm{s}}$ is the elevation-dependent shadow fading standard deviation following ITU-R P.681~\cite{itu_p681}.

\subsubsection{Doppler Model}
For a LEO satellite at altitude $d_{\mathrm{sat}}$ with orbital velocity $v_{\mathrm{sat}}$, the maximum Doppler shift is:
\begin{equation}\label{eq:fd_max}
  f_{\mathrm{d}} = \frac{v_{\mathrm{sat}} \cdot f_{\mathrm{c}}}{c},
\end{equation}
where $f_{\mathrm{c}}$ is the carrier frequency and $c$ is the speed of light. At $v_{\mathrm{sat}} = 7.22$~km/s and $f_{\mathrm{c}} = 2$~GHz, this yields $f_{\mathrm{d}} \approx 48.1$~kHz. The scatter paths exhibit a small Doppler spread around the LoS Doppler:
\begin{equation}\label{eq:doppler_scatter}
  \nu_i = f_{\mathrm{d}} + \delta_{\nu} \cos\!\left(\frac{2\pi(i-1)}{P}\right), \quad i = 1, \ldots, P,
\end{equation}
where $\delta_{\nu} = 0.05 f_{\mathrm{d}}$ represents the local scattering spread.

%% -----------------------------------------------------------------------
\subsection{DD-Domain Input-Output Relation}

When the DD grid $\mathbf{X}^{\mathrm{DD}}$ passes through the multipath channel, the received DD grid is given by a two-dimensional circular convolution:
\begin{equation}\label{eq:dd_io}
  Y^{\mathrm{DD}}[l,k] = \sum_{i=1}^{P} h_i\, X^{\mathrm{DD}}\!\left[(l - l_i)_N,\; (k - k_i)_M \right] + W[l,k],
\end{equation}
where $l_i = \mathrm{round}(\tau_i / \Delta\tau)$ and $k_i = \mathrm{round}(\nu_i / \Delta\nu)$ are the integer delay and Doppler bin indices, $(\cdot)_N$ denotes modulo-$N$ operation, and $W[l,k]$ is additive white Gaussian noise (AWGN).

\subsubsection{Fractional Delay-Doppler Spreading}
In practice, the channel parameters do not fall on exact integer bins. When fractional components exist, i.e., $\tau_i / \Delta\tau = l_i + \lambda_i$ and $\nu_i / \Delta\nu = k_i + \kappa_i$ with $\lambda_i, \kappa_i \in [0, 1)$, the received signal exhibits inter-bin interference modeled by the 2D Dirichlet kernel~\cite{raviteja2019}:
\begin{equation}\label{eq:dd_frac}
  Y^{\mathrm{DD}}[l,k] = \sum_{i=1}^{P} h_i \!\!\sum_{p=-N_{\mathrm{s}}}^{N_{\mathrm{s}}} \sum_{q=-N_{\mathrm{s}}}^{N_{\mathrm{s}}} \alpha_p^{(i)} \beta_q^{(i)}\, \Phi_i[l]\, X^{\mathrm{DD}}\!\left[(l\!-\!l_i\!-\!p)_N, (k\!-\!k_i\!-\!q)_M\right],
\end{equation}
where $N_{\mathrm{s}}$ is the spreading half-width, $\Phi_i[l] = e^{j2\pi l(k_i + \kappa_i)/(NM)}$ is the phase correction, and the delay and Doppler spreading coefficients are respectively given by:
\begin{align}
  \alpha_p^{(i)} &= \frac{1}{N} \cdot \frac{1 - e^{j2\pi(\lambda_i - p)}}{1 - e^{j2\pi(\lambda_i - p)/N}}, \label{eq:alpha}\\[4pt]
  \beta_q^{(i)} &= \frac{1}{M} \cdot \frac{1 - e^{j2\pi(\kappa_i - q)}}{1 - e^{j2\pi(\kappa_i - q)/M}}. \label{eq:beta}
\end{align}
These coefficients reduce to $\alpha_0 = \beta_0 = 1$ when the delay and Doppler are on-grid (i.e., $\lambda_i = \kappa_i = 0$), and \eqref{eq:dd_frac} degenerates to the integer model in~\eqref{eq:dd_io}.

%% -----------------------------------------------------------------------
\subsection{TF-Domain Equalization}

The 2D circular convolution structure in the DD domain translates to a diagonal TF-domain channel when delay and Doppler indices are integers. Specifically, the equivalent TF channel coefficient is:
\begin{equation}\label{eq:tf_otfs}
  H^{\mathrm{OTFS}}[n,m] = \sum_{i=1}^{P} h_i\, e^{j2\pi n l_i / N} \cdot e^{-j2\pi m k_i / M},
\end{equation}
which enables efficient one-tap MMSE equalization in the TF domain:
\begin{equation}\label{eq:mmse}
  \hat{X}^{\mathrm{TF}}[n,m] = \frac{H^{*}[n,m]}{|H[n,m]|^2 + \sigma_w^2 / \sigma_x^2} \cdot Y^{\mathrm{TF}}[n,m],
\end{equation}
where $\sigma_w^2$ and $\sigma_x^2$ are the noise and signal powers, respectively.

For comparison, the OFDM scheme processes each symbol in the time domain with per-sample Doppler phase rotation, which introduces inter-carrier interference (ICI). The OFDM one-tap equalizer uses:
\begin{equation}\label{eq:tf_ofdm}
  H^{\mathrm{OFDM}}[n,m] \!=\! \sum_{i=1}^{P} h_i \,\mathrm{sinc}(\nu_i^{\mathrm{r}} T) \, e^{j\pi \nu_i^{\mathrm{r}} T} e^{-j2\pi n l_i / N} e^{j2\pi \nu_i^{\mathrm{r}} m T_{\mathrm{s}}},
\end{equation}
where $\nu_i^{\mathrm{r}} = \nu_i - f_{\mathrm{d}}$ is the residual Doppler after receiver CFO correction, $T = 1/\Delta f$ is the useful OFDM symbol duration, and the $\mathrm{sinc}(\cdot)$ factor captures the amplitude loss due to ICI energy leakage~\cite{wang_ofdm_ici}. This ICI-induced loss is a fundamental disadvantage of OFDM in high-Doppler LEO channels.

%% -----------------------------------------------------------------------
\subsection{System Parameters}

Table~\ref{tab:params} summarizes the key system parameters adopted in this work, consistent with 5G NR NTN specifications.

\begin{table}[t]
\centering
\caption{System Parameters}
\label{tab:params}
\renewcommand{\arraystretch}{1.15}
\begin{tabular}{l|l}
\hline
\textbf{Parameter} & \textbf{Value} \\
\hline
Carrier frequency $f_{\mathrm{c}}$ & 2 GHz (S-band) \\
System bandwidth $B$ & 100 MHz \\
Subcarrier spacing $\Delta f$ & 120 kHz \\
Delay bins $N$ & 1024 \\
Doppler bins $M$ & 128 \\
Delay resolution $\Delta\tau$ & 8.14 ns \\
Doppler resolution $\Delta\nu$ & 876 Hz \\
CP ratio $\beta$ & 7\% \\
Modulation & 16-QAM \\
Channel coding & DVB-S.2 LDPC, rate 1/2 \\
LEO altitude & 600 km \\
Satellite velocity $v_{\mathrm{sat}}$ & 7.22 km/s \\
Elevation angle $\theta_{\mathrm{el}}$ & 50$^\circ$ \\
Channel model & NTN-TDL-C (3GPP TR 38.811) \\
K-factor & 12 dB (Suburban, 50$^\circ$) \\
Max delay spread & 90 ns ($\approx$11 bins) \\
Max Doppler shift $f_{\mathrm{d}}$ & $\approx$48.1 kHz \\
\hline
\end{tabular}
\end{table}

